\documentclass[12pt]{article}
\usepackage[T1]{fontenc}
\usepackage[utf8]{inputenc}
\usepackage{lmodern}
\usepackage{textcomp}
\usepackage{enumitem}
\usepackage{geometry}
\usepackage{graphicx}
\usepackage{amsmath, amssymb}
\geometry{margin=1in}
\begin{document}
\begin{center}
Political Science - Graduate Level Assessment 11 \
Total Marks: 40 \
Answer all questions.
\end{center}

\section*{Multiple Choice (10 marks)}
\begin{enumerate}[label=\arabic*.]
\item How many states in the international system are likely to have nuclear weapons right now?
\begin{enumerate}[label=(\alph*)]
    \item Fewer than 7
    \item Between 8 and 15
    \item Between 16 and 25
    \item More than 25
\end{enumerate}
\item What are the key elements of the Liberal approach to US foreign policy?
\begin{enumerate}[label=(\alph*)]
    \item Promotion of Democracy, free-trade and international institutions
    \item Alliances, diplomacy and protectionism
    \item The balance of power, self-sufficiency and prudence
    \item None of the above
\end{enumerate}
\item The concept of R$_{2}$P rests on three pillars, which of these is widely acclaimed to be the most important?
\begin{enumerate}[label=(\alph*)]
    \item The responsibility to retaliate.
    \item The responsibility to prevent.
    \item The responsibility to rebuild.
    \item The responsibility to react.
\end{enumerate}
\item In what sense might exceptionalism link isolationist and internationalist strategies?
\begin{enumerate}[label=(\alph*)]
    \item Both encourage world government
    \item Both focus on the decline of the American power
    \item It doesn't - the two are fundamentally opposed
    \item Both can be viewed as different means of achieving the same liberal ends
\end{enumerate}
\item The cooperative international organization of 185 countries designed to stabilize the exchange of currencies and the world economy is
\begin{enumerate}[label=(\alph*)]
    \item the World Bank.
    \item the United Nations.
    \item UNICEF.
    \item the International Monetary Fund.
\end{enumerate}
\end{enumerate}

\section*{True / False (10 marks)}
\textit{Answer True or False}
\begin{enumerate}[label=\arabic*.]
\setcounter{enumi}{5}
\item True or False: The costs of free trade are often concentrated among specific industries, while the benefits are widely dispersed across the economy.
\item True or False: All knowledge is a fixed and objective process that remains unaffected by social contexts.
\item True or False: Within Critical Security Studies, there is no specific referent object of security, emphasizing a broader understanding of security issues.
\item True or False: US foreign policy decision making can be constrained by the foreign policies of other states, international law, and intergovernmental organizations.
\item True or False: The grand strategy of 'primacy' refers to striving for American dominance in the international system.
\end{enumerate}

\section*{Long Form Response (20 marks)}
\textit{Answer each question with a clear and complete explanation.}
\begin{enumerate}[label=\arabic*.]
\setcounter{enumi}{10}
\item Discuss the objections realists have towards the influence of exceptionalism on American foreign policy, particularly concerning its impact on security, power, and national interests.
\item Discuss the distinctive phases in the development of security theory within security studies and analyze how these phases contribute to its separation from other related fields.
\end{enumerate}

\vfill
\noindent\textit{End of Paper}
\end{document}